\documentclass[../cours_2de.tex]{subfiles}
\begin{document}
\setcounter{chapter}{20}
\chapter{Probabilités finies et dénombrement}
\label{sec:ch21}
\begin{objectif}
Définir un modèle probabiliste fini ; calculer des événements composés ; distinguer loi supposée, calcul démontré et fréquence observée.
\end{objectif}
\section{Choisir un modèle}
Une expérience aléatoire est décrite par un univers fini $\Omega$ de résultats possibles, appelés issues.
Un événement est un sous-ensemble de $\Omega$.
Une loi de probabilité associe à chaque issue $\omega$ un nombre $p_\omega\geq0$ avec
$\sum_{\omega\in\Omega}p_\omega=1$.
On définit $P(A)=\sum_{\omega\in A}p_\omega$.
\begin{attention}
Le choix de la loi est une hypothèse du modèle, pas un théorème sur l'objet réel.
Un dé supposé équilibré conduit à six probabilités $1/6$. Une expérience réelle peut s'écarter de ce modèle.
\end{attention}
\begin{propriete}
$P(\varnothing)=0$, $P(\Omega)=1$, $0\leq P(A)\leq1$.
Si $A\subseteq B$, alors $P(A)\leq P(B)$.
Pour des événements disjoints, $P(A\cup B)=P(A)+P(B)$.
\end{propriete}
\begin{demonstration}
Les résultats découlent de sommes de nombres non négatifs.
Si $A\subseteq B$, les termes de $P(A)$ font partie de ceux de $P(B)$.
Si les événements sont disjoints, aucune issue n'est comptée deux fois dans la somme.
\end{demonstration}
\section{Contraire, réunion et intersection}
\begin{propriete}
$P(\overline A)=1-P(A)$ et
$P(A\cup B)=P(A)+P(B)-P(A\cap B)$.
\end{propriete}
\begin{demonstration}
$A$ et $\overline A$ partitionnent $\Omega$ : leurs probabilités somment à $1$.
Dans $P(A)+P(B)$, les issues de $A\cap B$ sont comptées deux fois ; les soustraire une fois donne exactement la somme des issues de la réunion.
\end{demonstration}
\begin{exemple}
Sur un dé équilibré, $A=\{2,4,6\}$ (pair), $B=\{4,5,6\}$ (au moins $4$).
Alors $P(A)=P(B)=1/2$, $P(A\cap B)=2/6=1/3$ et $P(A\cup B)=2/3$.
Le « ou » comprend les issues communes.
\end{exemple}
\section{Équiprobabilité et couples}
Si toutes les issues ont la même probabilité et si leur nombre est $N$, chacune vaut $1/N$, car leur somme vaut $1$.
Ainsi $P(A)=\operatorname{Card}(A)/N$.
Cette formule exige l'équiprobabilité des \emph{issues élémentaires}, pas celle de résultats regroupés.
\begin{exemple}[Deux dés]
Deux dés équilibrés lancés indépendamment donnent $36$ couples ordonnés équiprobables.
La somme $7$ correspond à six couples : $(1,6),(2,5),\ldots,(6,1)$ ; sa probabilité vaut $1/6$.
La somme $2$ ne correspond qu'à $(1,1)$, de probabilité $1/36$.
Les onze sommes possibles ne sont pas équiprobables.
\end{exemple}
Le produit cartésien rend l'ordre explicite. Pour un choix de $a$ possibilités suivi d'un choix de $b$ possibilités pour chacune, le nombre de couples est $ab$.
Un tableau ou un arbre de dénombrement permet de contrôler qu'aucune issue n'est oubliée ni comptée deux fois.
\section{Avec et sans remise}
\begin{exemple}
Une urne contient trois boules rouges distinctes et deux bleues distinctes.
Deux tirages avec remise donnent $5\times5=25$ couples équiprobables : neuf sont rouge--rouge, donc $P(RR)=9/25$.
Sans remise, il y a $5\times4=20$ couples ordonnés distincts, dont $3\times2=6$ rouge--rouge : $P(RR)=3/10$.
\end{exemple}
Le modèle est différent parce que le second tirage dépend du premier lorsqu'on ne remet pas la boule. Le chapitre suivant traduira ce raisonnement par des probabilités conditionnelles.
\section{Fréquence et probabilité}
La fréquence $f_n$ est observée dans $n$ répétitions ; la probabilité $p$ appartient au modèle.
Pour des répétitions indépendantes d'une même expérience, une version vulgarisée de la loi des grands nombres affirme que, pour un grand nombre de répétitions, les écarts importants entre $f_n$ et $p$ deviennent peu probables.
Elle ne dit ni $f_n=p$ exactement, ni que l'écart diminue à chaque nouveau tirage.
\begin{exemple}
Après dix piles consécutifs avec une pièce équilibrée et des lancers indépendants, le prochain lancer garde une probabilité $1/2$ de donner pile. Il n'existe pas d'obligation de « rattraper » les résultats précédents.
\end{exemple}
La simulation du chapitre 23 permet d'observer des fluctuations. Une justification quantitative facultative est donnée au chapitre 24 ; elle précise les hypothèses d'indépendance.
\utilite{Le calcul probabiliste est rigoureux une fois le modèle choisi. Savoir identifier les issues élémentaires est plus important que choisir immédiatement une formule.}
\begin{retenir}
Décrire l'univers et la loi, puis l'événement. Ne pas confondre « incompatible », « contraire » et « indépendant », ni fréquence expérimentale et probabilité.
\end{retenir}
\end{document}
